% META: {"id":"1SPE-DERLOCAL-EX-021","chapitre":"1SPE-DERIVATION-LOCAL","type_objet":"exercice","capacites":["1SPE-DERIVATION-LOCAL-C4"],"capacites_codes":["C4"],"methodes":["M4"],"parcours":2,"competences":["calculer","raisonner"],"duree_min":12,"mode_creation":"ex_nihilo","sources_inspiration":[],"fichier_tex":"chapitres/1SPE-DERIVATION-LOCAL/exercices/1SPE-DERLOCAL-EX-021.tex","corrige_tex":"chapitres/1SPE-DERIVATION-LOCAL/corriges/1SPE-DERLOCAL-CO-021.tex","status":"approved","origin":"generated","status_history":["generated","audited","accepted","approved"],"release_acceptance":"RELEASE_OWNER_BATCH_ACCEPTANCE","acceptance_closure_digest":"sha256:9b3ccf9a81c5520fbb7e03b7b2d3e7bf2057a02834b6d908bf3be5f49f3b553f","acceptance_receipt_digest":"sha256:4fad86f0282e0fa4e0419cca1ad6379ae7af5a280c04a4a90880f90a1c044242"}
% BEGIN-VERIFY
% from sympy import symbols, Eq, solve, Rational
% x = symbols('x')
% # f(x) = x^2 - 4x + 5, f'(x) = 2x - 4
% # a = 3 : f(3) = 9 - 12 + 5 = 2, f'(3) = 6 - 4 = 2
% # Tangente : y = 2(x - 3) + 2 = 2x - 4
% T = 2*(x - 3) + 2
% assert T.expand() == 2*x - 4
% assert T.subs(x, 3) == 2
% # Coupe Ox : 2x - 4 = 0 => x = 2
% sol = solve(Eq(T, 0), x)
% assert sol == [2]
% END-VERIFY
\begin{exercice}{1SPE-DERLOCAL-EX-021}{2}{12}
Soit $f$ la fonction définie sur $\mathbb{R}$ par $f(x) = x^2 - 4x + 5$. On admet que $f'(x) = 2x - 4$.

\begin{enumerate}
  \item Calculer $f(3)$ et $f'(3)$.
  \item Écrire l'équation de la tangente $T$ à la courbe de $f$ au point d'abscisse $3$.
  \item La tangente $T$ coupe-t-elle l'axe des abscisses ? Si oui, déterminer les coordonnées du point d'intersection.
\end{enumerate}
\end{exercice}
